\chapter*{第六部分导读：把课程项目做成可信交付}
\addcontentsline{toc}{chapter}{第六部分导读：把课程项目做成可信交付}

Part VI 是全书的 capstone 部分。到这里，学生已经学过计算环境、数据处理、机器学习、天文/物理案例、深度学习和现代 AI 工具。最后的问题不再是“还要学一个什么新模型”，而是：能否把前面所有能力组织成一个完整、可复查、可展示、可维护的小型科研项目。

Capstone 项目不是一份 notebook，也不是一份 PDF 报告，更不是一张好看的结果图。它至少应包含：

\begin{itemize}
    \item 清楚的问题定义和范围控制；
    \item 可复现的数据、代码和环境入口；
    \item baseline、模型、评估和错误分析；
    \item 图表、文字解释和限制说明；
    \item Git 提交记录、报告、展示材料和最终交付清单。
\end{itemize}

因此，Part VI 的章节看起来会更像项目管理、课程运行和交付规范。这不是偏离科研，而是把科研训练从“会做一次分析”推进到“能把分析交给别人复查和使用”。

\section*{Capstone 的核心产物}

一个合格的课程项目，至少应让助教或同学回答下面的问题：

\begin{enumerate}
    \item 这个项目想解决什么科学问题；
    \item 数据从哪里来，许可证和处理方式是什么；
    \item 运行入口在哪里，环境如何恢复；
    \item baseline 和主要模型分别是什么；
    \item 评估结果是否和错误案例一起解释；
    \item 最终结论有哪些限制；
    \item 如果下学期继续维护，应该从哪里开始。
\end{enumerate}

Part VI 的许多章节，就是围绕这些问题建立模板、rubric、handout、发布包、维护清单和交接机制。

\begin{protip}[交付不是最后一步，而是从第一天开始设计]
如果项目一开始没有范围、数据、环境和记录意识，最后很难靠临时整理变成可信交付。Capstone 的目标是让学生从项目第一周就知道最终需要交什么、为什么交、别人如何复查。
\end{protip}
